%==================================================================
\section{Discussion}
\label{sec:discussion}

\subsection{What the evaluation implies for practice}
\label{sec:practica}

The width that an adaptive conformal layer saves against a static split is not a measure of how much distribution shift the layer absorbed. To first order it is the conversion of the static split's coverage error into width, which any layer that targets nominal coverage performs: the constant-threshold oracle, which adapts nothing in time, reproduces \cf{mec_r2_pct}\% of the variation of the observed change across cells, in the real panel and, with known ground truth, in the simulation. The static split's coverage error in turn measures how far the calibration of the base predictive drifted between the calibration and the evaluation window. With the window held fixed, so that every base model faces the same shift, that drift spans \cf{oc_rango_transicion}~percentage points across base models in the transition window, and the base model explains \cf{anova_modelo_pct}\% of its variance against \cf{anova_ventana_pct}\% for the window. The simulation shows the same structure: a correctly specified predictive keeps the static split near nominal coverage under a tripling of magnitudes, a dispersion that is wrong by the same factor in every period is absorbed by the conformal quantile, and coverage drifts only when the base model fails to follow a change in the data.

These results suggest three checks before a gain is attributed to an adaptive calibration layer: report the coverage of the static split in the evaluation window; report the width of the constant-threshold oracle at the adaptive layer's coverage, so that the residual, which is what adaptation adds, is visible; and repeat the comparison on a second base model. All three are cheap, because a well-designed layer touches the base model only through its CDF and quantile function.

The ordering of effort follows. Replacing the reference hurdle by its early-stopped version, without touching the calibration layer, lowers the whole-test interval score of the static split by \cf{es_vs_ref_static_IS_pct}\%, and replacing it by the quantile model lowers it by \cf{qg_vs_ref_static_IS_pct}\%. The best calibration layer on the reference hurdle lowers it by \cf{ref_best_layer_IS_pct}\%. On a good base model the choice of layer matters little, and the ranking of layers itself depends on the base model: one-sided CQR is among the best layers on the quantile model and the worst on the early-stopped hurdle, because its correction is a single number of MWh pooled across plants of very different size (Section~\ref{sec:agnostic}).

For an operator or asset owner working with public data at this horizon, the evidence supports the following practice.
\begin{enumerate}
\item Invest first in the base predictive: a flexible distributional model, regularised or early-stopped, whose PIT histogram is checked on the calibration window and on recent data.
\item Wrap it in a static conformal layer, a PIT split or a one-sided CQR, recalibrated periodically.
\item If drift is expected, add ACI with a step selected out of sample and a lower bound on the miscoverage level, $\alpha_{\min}=0.005$, which removes infinite limits at a cost of at most half a point of coverage.
\item Check conditional coverage on regimes known when the forecast is issued, such as the trailing system curtailment, and report energy-weighted coverage next to plant-day coverage.
\item Do not use the transport-based scheme of Algorithm~\ref{alg:transporte}: once ACI is present the transport worsens the interval score, its tail shrinkage consumes most of the runtime for a negligible gain, and on outcome-defined high-curtailment days it leaves more curtailed energy uncovered than the static split.
\end{enumerate}

\subsection{The change in the record, and what we do not claim about it}
\label{sec:causas}

The distributional change documented in Section~\ref{sec:regime} is a ramp of roughly eighteen months between mid-2023 and late 2024. The dataset cannot identify its causes, because it records curtailment and not the assets or network states that absorb it. Several mechanisms overlap: growth of solar capacity behind a saturated north--south corridor, reinforcements of that corridor, changes in northern mining demand, the entry of battery storage, and the saturation of curtailment once the corridor is full through the solar hours. Mohamad et al.~\cite{mohamad2021} show that the effect of storage on curtailment depends on its position relative to the binding constraint, and Denholm and Mai~\cite{denholm2019} that it depends on duration as much as on power; we have neither the network model nor the charge-discharge records that such an analysis requires.

A second phenomenon appears at the turn of 2024 to 2025: a break in the rate of growth. The year-over-year growth of the monthly mean log-magnitude falls to $+0.09$ per year for northern solar from January 2025 and to $+0.01$ at system level from December 2024, while the northern solar fleet keeps expanding by 8 per cent through 2025. Storage entry is one candidate for that flattening and is of the right order, but the supporting energy balance is an upper bound that assumes one full daily cycle, unit efficiency and that every charged megawatt-hour would otherwise have been curtailed, and the same exercise for January to May 2026 does not reproduce it. We record storage as one hypothesis among several and make no causal claim. The supplementary material lists the storage systems of the SEN with their commercial operation dates, so that the hypothesis can be tested by anyone with network and dispatch data.

The calibration layers respond to only one of these two kinds of change. A displacement of the level of the predictive law appears as a displacement of the score distribution and is what the transport map and the online update are built to absorb. A break in the rate of growth leaves no displacement to correct and is invisible to coverage, so reliability under shifts of level carries no guarantee under changes of trend.

\subsection{The mid-2025 excursion}
\label{sec:excursion}

The coverage dip visible in Figure~\ref{fig:rolling} is concentrated in June 2025, the largest deseasonalised monthly anomaly of the post-calibration record. The Wasserstein-1 distance between May to July 2025 and August to December 2025 is 1.334, larger than the 1.147 of the transition contrast. Its signature is the reverse of the 2024 window: magnitudes rise while the occurrence probability falls. It appears in solar and in wind and across four regions, which rules out a single plant or a residual erratum, and the operator's report for that month records no storage system entering operation. A one-month excursion of this kind is absorbed only slowly by a continuously adapting scheme and is the case a change-point trigger would catch.

\subsection{Open directions}
\label{sec:direcciones}

Four directions follow from the results. A change-point detector could trigger recalibration at the onset of a shift, which would address excursions like that of June 2025; the offline machinery of Section~\ref{sec:regime} is a starting point. A scale-invariant version of CQR, for instance with the nonconformity score normalised by a predicted scale or computed on the log scale, would remove the dependence of the correction on the mix of plant sizes and on changes in level. Localised conformal prediction~\cite{guan2022}, which weights the calibration set by proximity to the test point, is a more principled route to conditional validity than the Mondrian partitions of \ref{app:extra}. And the evaluation can be extended to portfolios of plants that share a transmission corridor and whose curtailment is jointly determined.

%==================================================================
\section{Limitations}
\label{sec:limitations}

\begin{enumerate}
\item \emph{One system, one public data source.} All evidence comes from the SEN between 2022 and mid-2026. The mechanical part of the decomposition (Section~\ref{sec:mecanico}) follows from how a coverage-targeting layer works and should hold elsewhere; its size, the coefficients of the cross-cell relationship and the ranking of methods should not be transported to other systems without checking.
\item \emph{Two model families.} The nine base predictives are four variants of a lognormal hurdle and five capacity variants of a multi-quantile gradient-boosting model. The attribution of the coverage drift to the base model is established within this design; it is an empirical association for these families on this panel, and the intervals of Section~\ref{sec:mecanico} quantify sampling uncertainty over days, not uncertainty over model families.
\item \emph{The oracle is a benchmark.} The constant-threshold oracle uses the outcomes of each window and cannot be deployed. It separates the part of a width change that requires no adaptation from the rest; it does not provide an alternative method.
\item \emph{Dependence.} Exchangeability fails across the regime boundary and within it; the effective calibration sample is about 1{,}000 rows rather than 26{,}552 (Section~\ref{sec:guarantees}). Block, group and plant-level calibration (\ref{app:extra}) quantify this but do not repair it. Differences of one to three points of coverage are within the date-clustered standard error and are not interpreted.
\item \emph{High-curtailment days.} Coverage on days selected by the outcome is a description of exposure, not a test of calibration~\cite{lerch2017}; the ex-ante definition H3 is the legitimate test, and all definitions depend on a threshold, whose effect is reported in \ref{app:sens}.
\item \emph{A stylised cost.} The cost model is linear in the allowance and in the exceedance, with constant unit costs, and ignores prices, contracts and the timing of settlement. It ranks methods for three cost ratios; it does not value them in money.
\item \emph{A stylised simulation.} The simulated panel is generated from lognormal hurdles with common autoregressive shocks and linear ramps. It demonstrates the mechanism with known ground truth; its magnitudes are not estimates of the real ones.
\item \emph{Selection on a short window.} Base-model predictions start in January 2024, so the rolling-origin selection uses four months of inner validation. That is enough to reject steps that produce infinite limits, but a thin basis for fine distinctions among the gentler ones.
\item \emph{Point forecasting.} The comparison of Section~\ref{sec:punto} covers five model classes, twelve features, a seven-day horizon and mean absolute error. It says nothing about models with access to network topology, dispatch schedules, weather forecasts or storage state of charge, nor about criteria that weight large days.
\item \emph{What the dataset contains.} The dataset records curtailment only, with no network state, dispatch decision, price or storage operation, so the causes of the distributional change are not identified (Section~\ref{sec:causas}).
\end{enumerate}

%==================================================================
\section{Conclusion}
\label{sec:conclusion}

We have released an open, errata-corrected, plant-level dataset of renewable curtailment in Chile's National Electricity System, covering 53 continuous months and 300 plants and reconciling exactly with the operator's annual totals, and used it to evaluate conformal calibration of seven-day-ahead upper prediction limits under a documented change in the distribution of curtailment.

The main findings are four. First, the apparent benefit of an adaptive calibration layer depends on the base model: the scheme that narrows the reference hurdle's intervals by a third in the most divergent window leaves those of a multi-quantile model essentially unchanged. Second, to first order that benefit is the mechanical conversion of the static split's coverage error into width, which an oracle without any adaptation reproduces and a controlled simulation confirms; it measures the drift of the base predictive's calibration, not the size of the shift. Third, improving the base model dominates refining the calibration layer: early stopping alone narrows the hurdle's static limits by \cf{es_vs_ref_static_W_pct}\% at the same coverage, more than any calibration layer achieves on the reference hurdle. Fourth, on high-curtailment days defined by the outcome every method loses coverage and only those that over-cover on average stay at the nominal level, while on days identified from information available at the forecast origin the methods built on the quantile model stay close to nominal.

The transport-based adaptive pipeline is not the preferred operational method. On these data it under-covers on high-curtailment days (\cf{tr_alto_h1}\% against \cf{st_alto_h1}\% for the static split on the same base model), leaves more curtailed energy uncovered than the static split, and is dominated, at every cost ratio considered, by a static conformal wrapper (PIT split or one-sided CQR) around a multi-quantile predictive, which attains near-nominal coverage at less than half its width.

Three guarantees should be kept apart when such results are reported. The static split provides finite-sample marginal coverage under exchangeability, which this panel violates across the regime boundary and within it. The adaptive update provides long-run average coverage without distributional assumptions, and nothing at finite time. Empirical coverage in a dependent panel is a third thing: it is what we measure, and it varies across technologies, plant sizes, regions and high-curtailment days as reported. The dataset and the code that regenerates every table, figure and number are public.

%==================================================================
\section*{CRediT authorship contribution statement}
\textbf{Pablo Reyes Cerda:} Conceptualization, Data curation, Software, Validation, Investigation, Writing -- original draft. \textbf{Kerven Cea Morales:} Methodology, Formal analysis, Validation, Writing -- review and editing.

\section*{Declaration of competing interest}
Pablo Reyes Cerda is the founder of Synapta SpA and of CurtailmentIQ, a commercial forecasting service that uses data sources related to those described in this article. Kerven Cea Morales declares no competing interest.

\section*{Declaration of generative AI and AI-assisted technologies in the manuscript preparation process}
During the preparation of this work the authors used an AI-based assistant to support the drafting and language editing of the manuscript. After using this tool, the authors reviewed and edited the content as needed and take full responsibility for the content of the published article.

\section*{Data availability}
The dataset is openly available on Zenodo under a CC~BY~4.0 licence (DOI: 10.5281/zenodo.21198816, which resolves to the latest version of the deposit). The processing pipeline, the modelling code and the scripts that regenerate every table, figure and number in this article are available in the companion repository at \url{https://github.com/Pablosebastian-reyes/curtailmentiq-research}. The supplementary material contains a registry of the energy storage systems of the SEN with their commercial operation dates and sources, compiled from the public reports of the system operator, the Ministry of Energy and the National Energy Commission; it is not part of the released dataset.
